\begin{titlepage}
  \centering
  
  \scshape % Small Caps

  \vspace*{-2cm}
  
  %---------------------------------------------------------
  % Título
  
  \rule{0.97\textwidth}{1.6pt}\vspace*{-\baselineskip}\vspace*{2pt}
  \rule{0.97\textwidth}{0.4pt}
  
  \vspace{-0.2cm}
  
  \fontsize{40}{40}\selectfont Kifu\\
  \vspace{0.4cm}
  \fontsize{30}{30}\selectfont 棋譜\\
  \vspace{0.4cm}
  \fontsize{25}{25}\selectfont registro de partidas

  \vspace{-0.3cm}
  
  \rule{0.97\textwidth}{0.4pt}\vspace*{-\baselineskip}\vspace*{3pt}
  \rule{0.97\textwidth}{1.6pt}

  %---------------------------------------------------------

  \vspace*{0.1cm}

  %---------------------------------------------------------
  % Imagem
  
  % A partida histórica entre Shi Xiangxia (Branco) e Cheng Lanru (Preto) do século XVIII, "Nine Dragons Playing with a Pearl".

  \begin{figure}[h]
    \centering
    \hbox{
      \hspace{0.525cm}
      \begin{goban}[board dimension       = 15,
                    board size            = 19,
                    scale                 = 1,
                    outline line width    = 0.55mm,
                    horizontal clip start = 1,
                    horizontal clip end   = 19,
                    vertical clip start   = 1,
                    vertical clip end     = 19,
                    up to move            = 82]
        \parseSgfFile{\repoPath/sgf/kifu/nine_dragons_pearl.sgf}
      \end{goban}
    }
  \end{figure}

  %---------------------------------------------------------
  
  \vspace*{-0.15cm}

  %---------------------------------------------------------
  % Autor

  \fontsize{21}{21}\selectfont edição por\\
  \fontsize{28}{28}\selectfont Philippe Fanaro

  %---------------------------------------------------------
\end{titlepage}